% chapters/chap2.tex

\chapter{\texttt{\textbackslash includeonly} による選択コンパイル}

\texttt{\textbackslash includeonly\{ファイル1,ファイル2,...\}} をプリアンブルに書くと、
指定したファイルだけがコンパイルされます。

\section{選択コンパイルの利点}

大きな文書で作業するとき、全章をコンパイルすると時間がかかります。
\texttt{\textbackslash includeonly} を使えば、
\textbf{編集中の章だけ}を素早く確認できます。

\section{章番号・ページ番号の保持}

\texttt{\textbackslash includeonly} で一部の章だけコンパイルしても、
章番号・ページ番号・相互参照の番号は\textbf{全章分の情報が保たれます}。
これは \texttt{\textbackslash include} が各ファイルに対して
個別の \texttt{.aux} ファイルを生成しているためです。
